%=======================================================================
%  41277 Control Design -- Spring 2026
%  Assessment Task 5: Team Project Report (QUBE-Servo 2 rotary pendulum)
%
%  Build from this folder:  latexmk -pdf main.tex
%  Numbers quoted as \Macro{} come from results.tex, written by
%  make_figures.m together with every plot in figures/. Re-run that
%  script after any design change, then rebuild.
%  \todo{...} marks content only the team can supply (red in the PDF).
%=======================================================================
\documentclass[11pt,a4paper]{article}

\usepackage[T1]{fontenc}
\usepackage[utf8]{inputenc}
\usepackage[lining]{ebgaramond}
\usepackage[margin=2.2cm,top=3.0cm,headheight=42pt]{geometry}
\usepackage{amsmath,amssymb,bm}
\usepackage{mathtools}
\usepackage{graphicx}
\usepackage{float}
\usepackage{subcaption}
\usepackage{array}
\usepackage{booktabs}
\usepackage{fancyhdr}
\usepackage{enumitem}
\usepackage{xcolor}
\usepackage{listings}
\usepackage[colorlinks=true,linkcolor=black,urlcolor=blue!60!black,
            citecolor=black]{hyperref}

\setlength{\parindent}{0pt}
\setlength{\parskip}{6pt}
\allowdisplaybreaks
\raggedbottom

\newcommand{\todo}[1]{\textcolor{red}{\textbf{[TODO:} #1\textbf{]}}}
\lstset{language=Matlab,basicstyle=\ttfamily\footnotesize,frame=single,
        keywordstyle=\color{blue!70!black},commentstyle=\color{green!45!black},
        stringstyle=\color{purple!70!black},breaklines=true,columns=fullflexible}

\graphicspath{{figures/}}
\newcommand{\fig}[4][0.85]{%
  \begin{figure}[H]\centering\includegraphics[width=#1\linewidth]{#2}%
  \caption{#3}\label{#4}\end{figure}}

\newcommand{\KpTheta}{-3.994}
\newcommand{\KdTheta}{-1.724}
\newcommand{\KpAlpha}{42.90}
\newcommand{\KdAlpha}{2.384}
\newcommand{\PMcsA}{35.0}
\newcommand{\WcpcsA}{19.4}
\newcommand{\GMlowcsA}{0.53}
\newcommand{\GMupcsA}{Inf}
\newcommand{\MscsA}{1.66}
\newcommand{\GMlowcsAz}{0.53}
\newcommand{\GMupcsAz}{68}
\newcommand{\PMcsAz}{34.4}
\newcommand{\DMcsAz}{31}
\newcommand{\ZmaxcsA}{0.9944}
\newcommand{\OScsAL}{0.0}
\newcommand{\TscsAL}{1.64}
\newcommand{\UcsAL}{0.8}
\newcommand{\AlcsAL}{2.4}
\newcommand{\OScsBL}{0.6}
\newcommand{\TscsBL}{0.94}
\newcommand{\UcsBL}{1.6}
\newcommand{\AlcsBL}{4.6}
\newcommand{\OScsAN}{0.0}
\newcommand{\TscsAN}{1.64}
\newcommand{\UcsAN}{0.8}
\newcommand{\AlcsAN}{2.4}
\newcommand{\OScsBN}{0.6}
\newcommand{\TscsBN}{0.96}
\newcommand{\UcsBN}{1.6}
\newcommand{\AlcsBN}{4.6}
\newcommand{\GMlowcsB}{0.29}
\newcommand{\GMupcsB}{44.8}
\newcommand{\PMcsB}{67.7}
\newcommand{\FcsB}{[-19.4 110.195 -5.594 8.16]}
\newcommand{\FicsB}{-14.922}
\newcommand{\hwT}{171}
\newcommand{\hwAlphaStd}{0.97}
\newcommand{\hwAlphaMax}{12.1}
\newcommand{\hwThetaStd}{61.9}
\newcommand{\hwVrms}{0.79}
\newcommand{\hwGlitches}{2}
\newcommand{\hwRecover}{3.5}
\newcommand{\hwBumpPeak}{12.1}


\fancypagestyle{uts}{%
  \fancyhf{}
  \fancyhead[L]{\includegraphics[height=22pt]{../assets/logos/uts_logo.png}\quad
                \textbf{Task 5: Team Project Report}}
  \fancyhead[R]{\footnotesize\begin{tabular}{@{}r@{}}
                41277 Control Design\\ Spring 2026\\
                Coordinator: A/Prof Ricardo P. Aguilera
                \end{tabular}}
  \renewcommand{\headrulewidth}{0.4pt}
  \fancyfoot[C]{\thepage}
}
\pagestyle{uts}

\begin{document}

%=========================== COVER PAGE ================================
\thispagestyle{empty}
\includegraphics[height=30pt]{../assets/logos/uts_logo.png}\hfill
{\large\bfseries Faculty of Engineering and IT}

\vspace{10pt}
\fbox{%
\begin{minipage}{\dimexpr\textwidth-2\fboxsep-2\fboxrule\relax}
\vspace{4pt}
\centering{\Large\bfseries Assignment Cover Sheet}\par
\vspace{8pt}
\raggedright
\textbf{Subject:} 41277 Control Design\par
\vspace{4pt}
\textbf{Assignment Number:} Assessment Task 5
\hfill \textbf{Date Submitted:} \todo{date}\par
\vspace{4pt}
\textbf{Assignment Title:} Team Project Report: Frequency-Domain and
State-Feedback Control of the QUBE-Servo 2 Rotary Inverted Pendulum\par
\vspace{4pt}
\textbf{Student Name(s) and Number(s)}\hfill\textbf{Tutorial Group:} 02\par
\vspace{4pt}
\begin{tabular}{@{}p{4.2cm}p{3cm}p{5cm}@{}}
Jaylen Avtarovski & 24767939 & Contribution: \todo{\%}\\[4pt]
Rhyse Williams    & 24817532 & Contribution: \todo{\%}\\[4pt]
Elina Kim         & 24949885 & Contribution: \todo{\%}\\
\end{tabular}

\vspace{8pt}
\textbf{Declaration of Originality:}\par
{\small The work contained in this assignment, other than that specifically
attributed to another source, is that of the author(s). It is recognised
that, should this declaration be found to be false, disciplinary action
could be taken and the assignments of all students involved will be given
zero marks. In the statement below, I have indicated the extent to which I
have collaborated with other students, whom I have named.\par}

\vspace{6pt}
\textbf{Statement of Collaboration:}\par
\todo{The three team members above collaborated on all sections of this
report and on the design, simulation and hardware testing of both control
systems. Edit to describe how the team actually worked.}\par

\vspace{6pt}
\textbf{GenAI Declaration:}\par
\todo{Declare every generative-AI tool used and what it was used for
(for example drafting text, LaTeX formatting, MATLAB plotting scripts),
following the subject coordinator's advice. This must be accurate.}\par

\vspace{10pt}
\textbf{Signature} \quad \todo{initials} \hfill \textbf{Date} \todo{date}\par
\vspace{4pt}
\end{minipage}}

\vspace{12pt}
\textbf{Individual contribution levels.} As permitted by the assessment
criteria, the team states each member's contribution level, which scales
the individual report mark (individual mark $=$ group mark $\times$
contribution level).
\begin{table}[H]\centering
\begin{tabular}{lll}
\toprule
\textbf{Student}&\textbf{Contribution level}&\textbf{Main activities}\\
\midrule
Jaylen Avtarovski&\todo{\%}&\todo{}\\
Rhyse Williams&\todo{\%}&\todo{}\\
Elina Kim&\todo{\%}&\todo{}\\
\bottomrule
\end{tabular}
\end{table}

\clearpage
\tableofcontents
\clearpage

%=======================================================================
\section*{Acknowledgements}
\addcontentsline{toc}{section}{Acknowledgements}
\todo{Thank the people who helped: the teaching team and lab staff,
classmates, family, and so on.}

%=======================================================================
\section{Introduction}
The rotary inverted pendulum is a benchmark problem in control
engineering. A motor-driven horizontal arm must hold a freely hinged
pendulum upright, an open-loop unstable equilibrium, while also tracking a
commanded arm angle~\cite{furuta1992,astrom2000}. This project uses the
Quanser QUBE-Servo~2 with a Teensy~4.1 microcontroller running the
UTS ControLab Simulink template~\cite{qube_manual,hardware_template}.

The team designed, simulated and implemented two control systems that each
govern the whole system:
\begin{enumerate}[nosep]
\item \textbf{CS1, frequency domain.} Two transfer-function controllers,
  $C_\theta(s)$ for the arm and $C_\alpha(s)$ for the pendulum, plus a
  reference prefilter $F(s)$. They are designed in the $s$-domain and
  discretised to $C(z)$ for implementation.
\item \textbf{CS2, time domain.} A discrete-time state-feedback controller
  (SFC) with integral action and a Luenberger observer, designed from the
  discrete state-space model.
\end{enumerate}
Both controllers run inside a single MATLAB Function block at
$T_s=1$~ms. The same block is used by the linear simulation, the
nonlinear simulation and the code generated for the Teensy, so the code
that was simulated is the code that ran on the QUBE.

\textbf{Report structure.} Section~\ref{sec:lit} reviews the background,
Section~\ref{sec:objectives} states the control objectives,
Section~\ref{sec:model} develops and verifies the model, and
Section~\ref{sec:calib} covers sensor calibration.
Sections~\ref{sec:cs1} and~\ref{sec:cs2} present the design,
implementation and results of each control system.
Section~\ref{sec:compare} compares the two designs,
Section~\ref{sec:feedback} responds to the feedback from the final
demonstration, and Section~\ref{sec:conclusion} concludes.

%=======================================================================
\section{Literature review}\label{sec:lit}
Furuta \emph{et al.}~\cite{furuta1992} introduced the rotary pendulum as
a compact alternative to the cart--pole. It shares the cart--pole's
non-minimum-phase, underactuated dynamics but needs only a rotary joint
and no rail. Åström and Furuta~\cite{astrom2000} later separated the
problem into energy-based swing-up and linear stabilisation about the
upright equilibrium. This project addresses the second part: the
pendulum is raised by hand and the controller catches it near upright.

Classical frequency-domain design shapes the open-loop transfer function
so that the closed loop meets bandwidth and robustness targets, which are
expressed as gain and phase margins and the peak sensitivity
$M_s$~\cite{franklin_feedback,ogata}. For an open-loop unstable plant, the
Nyquist criterion requires the loop to encircle $-1$ once for each
unstable pole. The gain margin therefore becomes a \emph{band} of
admissible gains with both a lower and an upper limit.

State-space design instead places all closed-loop eigenvalues directly by
full-state feedback. That needs the velocities, which are not measured on
the QUBE, so they are reconstructed by an observer~\cite{aguilera2026,ogata}.
For digital implementation, the continuous designs are mapped to the
$z$-domain by Tustin or zero-order-hold equivalents, and continuous poles
map to $z=e^{sT_s}$~\cite{franklin_digital,mathworks_c2d}.

Inverted-pendulum stabilisation is the core problem in self-balancing
transporters, legged and humanoid robot balance, rocket attitude control
during vertical launch, and crane anti-sway control. This is why it is so
widely used as a teaching and research benchmark~\cite{ctms}.

%=======================================================================
\section{Control objectives}\label{sec:objectives}
Each control system must keep the pendulum upright ($\alpha\to0$) and move
the arm to a commanded angle $r$. Table~\ref{tab:specs} lists the targets.
The arm step used to assess every design is $r=45^\circ$. Its $\pm2\%$
settling band ($\pm0.9^\circ$) is about five encoder counts wide
($0.176^\circ$ per count). For a $5^\circ$ step the band would be smaller
than two counts, so the settling time could not be measured reliably.
\begin{table}[H]\centering
\caption{Closed-loop performance requirements.}\label{tab:specs}
\begin{tabular}{lll}
\toprule
\textbf{Quantity}&\textbf{Requirement}&\textbf{Reason}\\
\midrule
Arm overshoot $M_p$&$<15\%$&Project specification\\
Arm settling time $t_s$ (2\%)&$<5$~s&Project specification\\
Steady-state accuracy&\todo{confirm the specification}&Project specification\\
Motor voltage $|V_m|$&$\le10$~V&Amplifier limit~\cite{qube_manual}\\
Pendulum angle $|\alpha|$&small, about $<15^\circ$&Keeps the linear model valid\\
\bottomrule
\end{tabular}
\end{table}

%=======================================================================
\section{System modelling and model verification}\label{sec:model}
The model was derived in Demonstration Practical~2 and is reused here
unchanged, so both controllers are designed from the same plant.

\subsection{Nonlinear model and linearisation}

\textbf{Definitions and assumptions.}
The QUBE-Servo 2 rotary pendulum comprises a motor-driven horizontal arm
and a pendulum hinged at its tip. Encoders measure the two angles
\cite{qube_manual}. Let $\theta$ denote arm angle and $\alpha$ pendulum
angle from the upward vertical, with $\alpha=0$ upright. Use radians.
Let $\boldsymbol e_r$ point along the arm, $\boldsymbol e_t$ point in the
direction of increasing $\theta$ and $\boldsymbol e_z$ point upward.
Positive $\alpha$ tilts the pendulum towards $-\boldsymbol e_t$.
This explicitly fixes the coupling signs.

The arm and pendulum are approximated as rigid uniform slender rods.
The pendulum centre of mass is at $l_p=L_p/2$. Viscous damping represents
friction. Motor inductance, structural flexibility, nonlinear friction
and pendulum longitudinal-axis inertia are neglected. The motor is direct
drive. These are nominal modelling assumptions, not measurements of
the particular laboratory unit.
\begin{equation}
x=\begin{bmatrix}\theta&\alpha&\dot\theta&\dot\alpha\end{bmatrix}^{T},
\qquad u=V_m,\qquad y=\begin{bmatrix}\theta&\alpha\end{bmatrix}^{T}.
\end{equation}
Thus there are $n=4$ states, $m=1$ input and $p=2$ outputs. Velocities
are states even though only the angles are directly measured.

\textbf{Motor electrical dynamics.}
Kirchhoff's voltage law and the motor relations give
\begin{align}
V_m&=R_mi_m+L_m\frac{di_m}{dt}+e_b,\\
e_b&=k_m\dot\theta,\qquad \tau_m=k_ti_m.
\end{align}
Neglecting inductance,
\begin{equation}\label{eq:motor}
i_m=\frac{V_m-k_m\dot\theta}{R_m},\qquad
\boxed{\tau_m=\frac{k_t}{R_m}V_m-\frac{k_tk_m}{R_m}\dot\theta.}
\end{equation}
Voltage generates torque and back emf reduces the current as speed rises.

\textbf{Inertias and energy.}
For a uniform rod, $dm=(m/L)\,dr$. The relevant inertias are
\begin{align}
J_{\mathrm{arm}}&=\int_0^{L_r}r^2\frac{m_r}{L_r}\,dr
=\frac13m_rL_r^2,\qquad J_r=J_{\mathrm{arm}}+J_m+J_h,\\
J_p&=\int_{-L_p/2}^{L_p/2}r^2\frac{m_p}{L_p}\,dr
=\frac1{12}m_pL_p^2.
\end{align}
Here $J_r$ includes the arm, rotor and hub. The arm rod approximation
excludes the separately added rotor and hub. No inertia-disc module is
included. $J_p$ is about the pendulum centre of mass; the hinge inertia is
$J_p+m_pl_p^2=\frac13m_pL_p^2$ by the parallel-axis theorem.

The pendulum centre-of-mass position is
\begin{equation}
\boldsymbol r_p=L_r\boldsymbol e_r-l_p\sin\alpha\,\boldsymbol e_t
+l_p\cos\alpha\,\boldsymbol e_z.
\end{equation}
Using $\dot{\boldsymbol e}_r=\dot\theta\boldsymbol e_t$ and
$\dot{\boldsymbol e}_t=-\dot\theta\boldsymbol e_r$ gives
\begin{align}
\dot{\boldsymbol r}_p&=l_p\sin\alpha\,\dot\theta\,\boldsymbol e_r
+(L_r\dot\theta-l_p\cos\alpha\,\dot\alpha)\boldsymbol e_t
-l_p\sin\alpha\,\dot\alpha\,\boldsymbol e_z,\\
\|\dot{\boldsymbol r}_p\|^2&=(L_r^2+l_p^2\sin^2\alpha)\dot\theta^2
+l_p^2\dot\alpha^2-2L_rl_p\cos\alpha\,\dot\theta\dot\alpha.
\end{align}
Neglecting longitudinal-axis inertia, pendulum rotation about its centre
of mass contributes
\[\frac12J_p(\dot\alpha^2+\sin^2\alpha\,\dot\theta^2).\]
Therefore the total kinetic and potential energies are
\begin{align}
T&=\frac12J_r\dot\theta^2+\frac12m_p\|\dot{\boldsymbol r}_p\|^2
+\frac12J_p(\dot\alpha^2+\sin^2\alpha\,\dot\theta^2),\\
U&=m_pgl_p\cos\alpha.
\end{align}
Potential energy is referenced to the hinge height and is maximal upright.

Define
\begin{equation}\label{eq:constants}
a=J_r+m_pL_r^2,\quad b=J_p+m_pl_p^2,\quad
c=m_pL_rl_p,\quad h=m_pgl_p.
\end{equation}
The Lagrangian is
\begin{equation}\label{eq:lagrangian}
\mathcal L=T-U=
\frac12(a+b\sin^2\alpha)\dot\theta^2+\frac12b\dot\alpha^2
-c\cos\alpha\,\dot\theta\dot\alpha-h\cos\alpha.
\end{equation}

\textbf{Euler--Lagrange equations.}
For $q\in\{\theta,\alpha\}$,
\begin{equation}
\frac{d}{dt}\left(\frac{\partial\mathcal L}{\partial\dot q}\right)
-\frac{\partial\mathcal L}{\partial q}=Q_q,\qquad
Q_\theta=\tau_m-D_r\dot\theta,\quad Q_\alpha=-D_p\dot\alpha.
\end{equation}
The arm derivatives are
\begin{align}
\frac{\partial\mathcal L}{\partial\theta}&=0,\\
\frac{\partial\mathcal L}{\partial\dot\theta}
&=(a+b\sin^2\alpha)\dot\theta-c\cos\alpha\,\dot\alpha,\\
\frac{d}{dt}\left(\frac{\partial\mathcal L}{\partial\dot\theta}\right)
&=(a+b\sin^2\alpha)\ddot\theta
+2b\sin\alpha\cos\alpha\,\dot\alpha\dot\theta
-c\cos\alpha\,\ddot\alpha+c\sin\alpha\,\dot\alpha^2.
\end{align}
The pendulum derivatives are
\begin{align}
\frac{\partial\mathcal L}{\partial\dot\alpha}
&=b\dot\alpha-c\cos\alpha\,\dot\theta,\\
\frac{d}{dt}\left(\frac{\partial\mathcal L}{\partial\dot\alpha}\right)
&=b\ddot\alpha-c\cos\alpha\,\ddot\theta
+c\sin\alpha\,\dot\alpha\dot\theta,\\
\frac{\partial\mathcal L}{\partial\alpha}
&=b\sin\alpha\cos\alpha\,\dot\theta^2
+c\sin\alpha\,\dot\theta\dot\alpha+h\sin\alpha.
\end{align}
Define $k_v=k_t/R_m$ and $d=D_r+k_tk_m/R_m$. The resulting
nonlinear mechanical equations, including motor torque, are
\begin{samepage}
\begin{subequations}\label{eq:eom}
\begin{align}
(a+b\sin^2\alpha)\ddot\theta-c\cos\alpha\,\ddot\alpha
+2b\sin\alpha\cos\alpha\,\dot\alpha\dot\theta
+c\sin\alpha\,\dot\alpha^2&=k_vu-d\dot\theta,\\
-c\cos\alpha\,\ddot\theta+b\ddot\alpha
-b\sin\alpha\cos\alpha\,\dot\theta^2-h\sin\alpha&=-D_p\dot\alpha.
\end{align}
\end{subequations}
\end{samepage}
Fixing the arm and neglecting friction gives $b\ddot\alpha=h\sin\alpha$.
A small positive tilt therefore accelerates away from upright, which
provides a physical check on the gravity sign.

\textbf{Nonlinear state-space representation.}
Define the torque expressions
\begin{align}
r_1(x,u)&=k_vu-dx_3-2b\sin x_2\cos x_2\,x_4x_3-c\sin x_2\,x_4^2,\\
r_2(x,u)&=b\sin x_2\cos x_2\,x_3^2+h\sin x_2-D_px_4.
\end{align}
The equations of motion become
\begin{equation}
\underbrace{\begin{bmatrix}
a+b\sin^2\alpha&-c\cos\alpha\\-c\cos\alpha&b
\end{bmatrix}}_{M(\alpha)}
\begin{bmatrix}\ddot\theta\\\ddot\alpha\end{bmatrix}
=\begin{bmatrix}r_1\\r_2\end{bmatrix}.
\end{equation}
The mass-matrix determinant and inverse are
\begin{align}
\Delta(\alpha)&=b(a+b\sin^2\alpha)-c^2\cos^2\alpha,\\
M^{-1}(\alpha)&=\frac1{\Delta(\alpha)}
\begin{bmatrix}b&c\cos\alpha\\c\cos\alpha&a+b\sin^2\alpha\end{bmatrix}.
\end{align}
Hence
\begin{equation}\label{eq:nonlinear_ss}
\boxed{\dot x=f(x,u)=
\begin{bmatrix}
x_3\\x_4\\
(br_1+c\cos x_2\,r_2)/\Delta(x_2)\\
(c\cos x_2\,r_1+(a+b\sin^2x_2)r_2)/\Delta(x_2)
\end{bmatrix}.}
\end{equation}
For the parameters in Section~\ref{sec:model}, $ab-c^2>0$ and
$\Delta(\alpha)=(ab-c^2)+(b^2+c^2)\sin^2\alpha>0$, so this idealised
model is nonsingular for every pendulum angle.

\textbf{Equilibrium and Jacobian linearisation.}
An upright stationary equilibrium is
\begin{equation}
x_e=\begin{bmatrix}\theta_e&0&0&0\end{bmatrix}^{T},\qquad u_e=0.
\end{equation}
Any fixed $\theta_e$ is an ideal equilibrium, since no restoring torque
depends on arm angle. Choose $\theta_e=0$ for the numerical model.
Define $\delta x=x-x_e$, $\delta u=u-u_e$ and $\delta y=y-y_e$.
The first-order Taylor expansion is
\begin{equation}\label{eq:ct_ss}
\delta\dot x=A_c\delta x+B_c\delta u,\qquad
\delta y=C_c\delta x+D_c\delta u,
\end{equation}
with Jacobians
\begin{equation}
A_c=\left.\frac{\partial f}{\partial x}\right|_e,\qquad
B_c=\left.\frac{\partial f}{\partial u}\right|_e.
\end{equation}
At equilibrium, $r_1=r_2=0$ and $\Delta_0=ab-c^2$. Also,
\begin{equation}
\left.\frac{\partial(r_1,r_2)}{\partial x}\right|_e
=\begin{bmatrix}0&0&-d&0\\0&h&0&-D_p\end{bmatrix},\qquad
\left.\frac{\partial(r_1,r_2)}{\partial u}\right|_e
=\begin{bmatrix}k_v\\0\end{bmatrix}.
\end{equation}
Derivatives of $M^{-1}$ multiply the zero equilibrium torque vector.
Multiplying the above derivatives by
$M^{-1}(0)=\Delta_0^{-1}\begin{bmatrix}b&c\\c&a\end{bmatrix}$ gives
\begin{equation}\label{eq:Ac_symbolic}
\boxed{A_c=
\begin{bmatrix}
0&0&1&0\\0&0&0&1\\
0&ch/\Delta_0&-bd/\Delta_0&-cD_p/\Delta_0\\
0&ah/\Delta_0&-cd/\Delta_0&-aD_p/\Delta_0
\end{bmatrix},\quad
B_c=\begin{bmatrix}0\\0\\bk_v/\Delta_0\\ck_v/\Delta_0\end{bmatrix}.}
\end{equation}
Since the outputs are the two angles with no direct voltage feedthrough,
\begin{equation}\label{eq:CD}
\boxed{C_c=\begin{bmatrix}1&0&0&0\\0&1&0&0\end{bmatrix},
\qquad D_c=\begin{bmatrix}0\\0\end{bmatrix}.}
\end{equation}
This linearisation describes small deviations about upright; large-angle
motion requires the nonlinear model in~\eqref{eq:nonlinear_ss}.

\subsection{System parameters}

Table~\ref{tab:params} lists the manufacturer-supplied parameters relevant
to the pendulum configuration. The inertia-disc parameters are excluded.
Motor inductance is listed to justify its neglect, rather than being
retained as a model state.
\begin{table}[H]\centering
\caption{Parameters from Table 2.2 of the QUBE-Servo 2 user manual~\cite{qube_manual}.}
\label{tab:params}
\begin{tabular}{lll}
\toprule
\textbf{Symbol}&\textbf{Description}&\textbf{Value}\\
\midrule
$R_m$&Motor resistance&$8.4~\Omega$\\
$k_t$&Torque constant&$0.042~\mathrm{N\,m/A}$\\
$k_m$&Back-emf constant&$0.042~\mathrm{V\,s/rad}$\\
$J_m$&Motor rotor inertia&$4.0\times10^{-6}~\mathrm{kg\,m^2}$\\
$J_h$&Attachment hub inertia&$0.6\times10^{-6}~\mathrm{kg\,m^2}$\\
$L_m$&Inductance (neglected)&$1.16\times10^{-3}~\mathrm H$\\
$m_r$&Arm mass&$0.095~\mathrm{kg}$\\
$L_r$&Arm length, motor pivot to tip&$0.085~\mathrm m$\\
$m_p$&Pendulum mass&$0.024~\mathrm{kg}$\\
$L_p$&Pendulum full length&$0.129~\mathrm m$\\
&Encoder resolution (quadrature)&$2\pi/2048~\mathrm{rad/count}$\\
&Recommended amplifier output&$\pm10~\mathrm V$\\
\bottomrule
\end{tabular}
\end{table}
The motor electrical time constant is
\begin{equation}
\tau_{\mathrm{elec}}=\frac{L_m}{R_m}
=\frac{1.16\times10^{-3}}{8.4}
=1.3810\times10^{-4}~\mathrm s.
\end{equation}
This is short compared with the mechanical time scales, motivating the
quasi-static current approximation. Retaining inductance would introduce
motor current as a fifth state.

The damping values in Table~\ref{tab:additional} are adopted from the
manufacturer's student workbook, \emph{Pendulum State-Space Model}, in-lab
exercise 6~\cite{qube_courseware}. That source reports them as experimentally
obtained example values; they were not identified from this group's
specific laboratory unit. No new experimental identification is claimed.
\begin{table}[H]\centering
\caption{Additional constants and their provenance.}
\label{tab:additional}
\begin{tabular}{llll}
\toprule
\textbf{Symbol}&\textbf{Description}&\textbf{Value}&\textbf{Basis}\\
\midrule
$D_r$&Arm damping&$0.0015~\mathrm{N\,m\,s/rad}$&Workbook example\\
$D_p$&Pendulum damping&$0.0005~\mathrm{N\,m\,s/rad}$&Workbook example\\
$g$&Gravity&$9.81~\mathrm{m/s^2}$&Standard approximation\\
\bottomrule
\end{tabular}
\end{table}

\begin{samepage}
The derived dimensions and inertias are
\begin{align}
l_p&=L_p/2=0.0645~\mathrm m,\\
J_{\mathrm{arm}}&=\tfrac13(0.095)(0.085)^2
=2.2879167\times10^{-4}~\mathrm{kg\,m^2},\\
J_r&=J_{\mathrm{arm}}+4.0\times10^{-6}+0.6\times10^{-6}
=2.3339167\times10^{-4}~\mathrm{kg\,m^2},\\
J_p&=\tfrac1{12}(0.024)(0.129)^2
=3.3282\times10^{-5}~\mathrm{kg\,m^2}.
\end{align}
\end{samepage}
The constants used in the matrices are
\begin{align}
a&=4.0679167\times10^{-4}~\mathrm{kg\,m^2},&
b&=1.33128\times10^{-4}~\mathrm{kg\,m^2},\\
c&=1.3158\times10^{-4}~\mathrm{kg\,m^2},&
h&=0.01518588~\mathrm{N\,m},\\
d&=0.00171~\mathrm{N\,m\,s/rad},&
k_v&=0.005~\mathrm{N\,m/V},\\
\Delta_0&=3.68420646\times10^{-8}~\mathrm{kg^2\,m^4}.&&
\end{align}
Substituting into~\eqref{eq:Ac_symbolic} gives
\begin{equation}\label{eq:Ac_numeric}
A_c\simeq
\begin{bmatrix}
0&0&1&0\\0&0&0&1\\
0&54.2357795&-6.17904785&-1.78573054\\
0&167.674898&-6.10719845&-5.52075014
\end{bmatrix},\quad
B_c\simeq\begin{bmatrix}0\\0\\18.0673914\\17.8573054\end{bmatrix}.
\end{equation}
The matrices $C_c,D_c$ remain exactly those in~\eqref{eq:CD}.
These are nominal simulation parameters. Actual mass distribution, damping,
encoder calibration and unmodelled friction may produce differences from
hardware measurements.

\subsection{Discretisation and sampling time}

Choose the sampling time
\begin{equation}
\boxed{T_s=0.001~\mathrm s,\qquad f_s=1/T_s=1000~\mathrm{Hz}.}
\end{equation}
This matches the value in the supplied ControLab template's
\texttt{Hardware\_settings.m}~\cite{hardware_template}. It is also short
relative to the mechanical dynamics, since
$T_s\max_i|\lambda_i(A_c)|\simeq0.01763\ll1$ for this model. Hardware execution timing
and achieved closed-loop performance must still be checked during
implementation.

Zero-order hold (ZOH) assumes that motor voltage is held constant between
controller updates. The exact discretisation of the linearised model
under this assumption is
\begin{equation}\label{eq:zoh}
A=e^{A_cT_s},\qquad
B=\int_0^{T_s}e^{A_c\tau}B_c\,d\tau,\qquad C=C_c,\quad D=D_c.
\end{equation}
The MATLAB function \texttt{c2d} evaluates these expressions without
requiring an inverse of $A_c$. This matters because the arm-position mode
makes $A_c$ singular.
The discrete-time model is
\begin{subequations}\label{eq:dt_ss}
\begin{align}
\delta x[k+1]&=A\delta x[k]+B\delta u[k],\\
\delta y[k]&=C\delta x[k]+D\delta u[k].
\end{align}
\end{subequations}
At the chosen zero equilibrium the deviation variables equal the
original variables. The calculated matrices are
\begin{equation}\label{eq:A_discrete}
A\simeq
\begin{bmatrix}
1&2.70127657\!\times\!10^{-5}&9.96918625\!\times\!10^{-4}&-8.80391044\!\times\!10^{-7}\\
0&1.00008362958&-3.04176177\!\times\!10^{-6}&9.97274363\!\times\!10^{-4}\\
0&0.05392103923&0.99384536883&-0.00174835485\\
0&0.16705290505&-0.00607175726&0.99458335877
\end{bmatrix},
\end{equation}
\begin{equation}
B\simeq\begin{bmatrix}
9.00986699\times10^{-6}\\8.89404027\times10^{-6}\\
0.0179959975588\\0.0177536761881
\end{bmatrix},\quad
C=\begin{bmatrix}1&0&0&0\\0&1&0&0\end{bmatrix},\quad
D=\begin{bmatrix}0\\0\end{bmatrix}.
\end{equation}
The dimensions are $A\in\mathbb R^{4\times4}$,
$B\in\mathbb R^{4\times1}$, $C\in\mathbb R^{2\times4}$ and
$D\in\mathbb R^{2\times1}$. Displayed results are rounded;
full precision must be retained in MATLAB for controller and observer
design.



\subsection{Transfer functions for frequency-domain design}
From $G(s)=C_c(sI-A_c)^{-1}B_c$ (MATLAB \texttt{tf(ss(Ac,Bc,Cc,Dc))}),
both outputs share the denominator
\begin{equation}\label{eq:Dplant}
D(s)=s\,(s^3+d_3s^2+d_2s+d_1)=s^4+11.70s^3-144.47s^2-704.84s,
\end{equation}
where $d_3=(aD_p+bd)/\Delta_0$, $d_2=(dD_p-ah)/\Delta_0$ and
$d_1=-dh/\Delta_0$. The two transfer functions are
\begin{align}
G_\theta(s)=\frac{\theta(s)}{V_m(s)}&=
\frac{18.07s^2+67.86s-2060.9}{D(s)}
=\frac{18.07\,(s-8.97)(s+12.72)}{s\,(s-9.95)(s+4.02)(s+17.6)},\label{eq:Gth}\\
G_\alpha(s)=\frac{\alpha(s)}{V_m(s)}&=\frac{17.86\,s^2}{D(s)}
=\frac{17.86\,s}{(s-9.95)(s+4.02)(s+17.6)}.\label{eq:Gal}
\end{align}
These transfer functions show three properties that drive the design:
\begin{itemize}[nosep]
\item \textbf{Open-loop unstable.} The pole at $s=+9.95$~s$^{-1}$ is the
  falling pendulum (time constant about 0.1~s), and the pole at $s=0$
  reflects that the arm has no restoring force. A stabilising
  pendulum loop is mandatory.
\item \textbf{Non-minimum phase arm.} $G_\theta$ has a right-half-plane
  zero at $s=+8.97$. To move the arm positively, the base must first
  swing the wrong way to tilt the pendulum. This limits the achievable arm
  bandwidth and causes the initial undershoot seen in the responses.
\item \textbf{Zero at $s=0$ in $G_\alpha$.} A constant voltage cannot hold
  a steady tilt, so the pendulum loop alone cannot fix the arm position.
  An outer arm loop is required.
\end{itemize}

\subsection{Model verification}\label{sec:verify}
The linear model was verified against the nonlinear model by simulating
both with the same controller code. For the $45^\circ$ arm step, the
pendulum stays within $\AlcsAN^\circ$ (CS1) and $\AlcsBN^\circ$ (CS2) of
upright. The linear and nonlinear responses overlap within plotting
accuracy (Figures~\ref{fig:cs1sim} and~\ref{fig:cs2sim}), which confirms
that the small-angle assumption holds over the operating range.

Model signs were checked on the hardware. With the controller disabled and
a small positive voltage applied, $\theta$ increased as the model
predicts. The pendulum encoder measured $\alpha$ in the opposite direction
to the model's convention, so a gain of $-1$ was placed before the
controller (Section~\ref{sec:calib}). After this correction the CS1
controller balanced the pendulum on its first attempt, which independently
confirms the sign structure of~\eqref{eq:Gth}--\eqref{eq:Gal}.
\todo{If the team ran an open-loop or system-identification test, add a
measured-versus-model plot here. The DP2 damping values are textbook
estimates, not identified from this QUBE.}

%=======================================================================
\section{Calibration of transducers and sensors}\label{sec:calib}
Both angles are measured by quadrature encoders with 2048 counts per
revolution, a resolution of $2\pi/2048=3.07$~mrad ($0.176^\circ$) per
count~\cite{qube_manual}. The calibration steps were:
\begin{enumerate}[nosep]
\item \textbf{Zero of $\alpha$.} The pendulum encoder counts from the angle
  at power-up. The pendulum is therefore started hanging down, where the
  template measures $\alpha=-\pi$, and the template's offset makes upright
  $\alpha=0$, as the model requires. Each run starts by checking that the
  hanging reading is $-180^\circ$.
\item \textbf{Direction of $\alpha$.} A gain of $-1$ makes a positive
  tilt match the model's $-\bm e_t$ convention (Section~\ref{sec:verify}).
\item \textbf{Direction and zero of $\theta$.} The arm angle is zero at
  power-up. A positive voltage increases the measured $\theta$, matching
  the model, so no correction was needed.
\item \textbf{Motor voltage.} The commanded $V_m$ is saturated to
  $\pm10$~V inside the controller block, so the observer and the logged data
  always see the voltage the amplifier actually applies.
\item \textbf{Encoder read errors.} The logged arm angle contains
  occasional single-sample spikes of the order of $10^5$ degrees (\hwGlitches{} in
  run~04). These are SPI read errors, not motion, and they are removed
  before plotting by holding the previous sample whenever $\theta$ jumps
  by more than $90^\circ$ in 1~ms.
\end{enumerate}
\todo{Describe any further calibration the team did, for example
checking the encoder scale against a protractor at $\pm90^\circ$.}

%=======================================================================
\section{Control System 1: frequency-domain design}\label{sec:cs1}
\subsection{Control structure and why it was chosen}
Figure~\ref{fig:cs1block} shows the structure. The plant has one input and
two outputs, so a single SISO loop cannot both stabilise $\alpha$ and fix
$\theta$ (Section~\ref{sec:model}). Two controllers act on one voltage:
\begin{equation}\label{eq:cs1law}
V_m=\operatorname{sat}_{\pm10}\Big[C_\theta(s)\big(F(s)r-\theta\big)
-C_\alpha(s)\,\alpha\Big].
\end{equation}
\begin{figure}[H]\centering
\setlength{\unitlength}{1mm}
\begin{picture}(150,46)
  \put(0,36){\vector(1,0){12}}\put(2,38){$r$}
  \put(12,31){\framebox(14,10){$F(s)$}}
  \put(26,36){\vector(1,0){10}}
  \put(38,36){\circle{4}}\put(34,39){\small$+$}\put(39,31){\small$-$}
  \put(40,36){\vector(1,0){8}}
  \put(48,31){\framebox(16,10){$C_\theta(s)$}}
  \put(64,36){\line(1,0){8}}\put(74,26){\circle{4}}
  \put(72,36){\line(1,0){2}}\put(74,36){\vector(0,-1){8}}\put(76,31){\small$+$}
  \put(52,8){\framebox(16,10){$C_\alpha(s)$}}
  \put(68,13){\line(1,0){6}}\put(74,13){\vector(0,1){11}}\put(76,17){\small$-$}
  \put(76,26){\vector(1,0){10}}\put(78,28){\small$V_m$}
  \put(86,19){\framebox(26,14){\shortstack{QUBE\\$G_\theta,G_\alpha$}}}
  \put(112,30){\line(1,0){20}}\put(124,32){$\theta$}
  \put(132,30){\line(0,1){14}}\put(132,44){\line(-1,0){94}}\put(38,44){\vector(0,-1){6}}
  \put(112,22){\line(1,0){20}}\put(124,17){$\alpha$}
  \put(132,22){\line(0,-1){18}}\put(132,4){\line(-1,0){72}}\put(60,4){\vector(0,1){4}}
  \put(46,0){\dashbox{1}(72,46){}}
\end{picture}
\caption{CS1 structure. Everything inside the dashed box runs in one
MATLAB Function block at $T_s=1$~ms. The saturation is omitted for
clarity.}\label{fig:cs1block}
\end{figure}

\textbf{Controller type.} Both $C_\theta$ and $C_\alpha$ are filtered PD
controllers:
\begin{equation}\label{eq:pd}
C_i(s)=K_{p,i}+K_{d,i}\frac{Ns}{s+N},\qquad i\in\{\theta,\alpha\},\qquad
N=60~\mathrm{rad/s}.
\end{equation}
\begin{itemize}[nosep]
\item \emph{Derivative action} adds the phase lead needed to stabilise the
  unstable pole at $s=+9.95$. A proportional-only pendulum loop cannot do
  this, because the gravity term makes the loop phase too negative near
  crossover.
\item \emph{The first-order filter} at $N=60$~rad/s limits the gain on
  encoder quantisation noise. An ideal derivative of a signal that changes
  in $0.176^\circ$ steps every few milliseconds would drive the voltage
  between its limits. $N$ is placed about $3\times$ above the target
  crossover of about 20~rad/s, so the filter costs only about $18^\circ$ of
  phase there.
\item \emph{No integrator} is needed. The plant already contains the
  free integrator $1/s$ from arm angle to voltage, so the arm loop is type~1
  and has zero steady-state error to a step without friction.
\end{itemize}

\subsection{Design method}\label{sec:cs1design}
\textbf{Step 1: target closed-loop poles from the specifications.} The
dominant pair is chosen from the second-order relations
$M_p=e^{-\pi\zeta/\sqrt{1-\zeta^2}}$ and $t_s\approx4/(\zeta\omega_n)$:
\begin{equation}
\zeta=0.7\Rightarrow M_p=4.6\%<15\%,\qquad
\omega_n=8~\mathrm{rad/s}\Rightarrow t_s\approx\frac{4}{0.7\times8}=0.71~\mathrm s<5~\mathrm s.
\end{equation}
This leaves margin for the RHP zero, friction and quantisation, which all
slow and degrade the real response. The bandwidth was \emph{not} set by
the settling-time limit alone. An earlier, slower design ($\omega_n=2.5$,
$\zeta=0.6$) met the specifications in simulation but showed a persistent
arm oscillation once encoder quantisation was modelled. Its pendulum loop
had too little gain to reject the quantised measurement. The faster pole
set raises the pendulum-loop gain while keeping the voltage small.

Two further poles are placed at $-2.4\omega_n=-19.2$ and
$-3.2\omega_n=-25.6$, at least $3\times$ faster than the dominant pair, so
they add little to the transient.

\textbf{Step 2: characteristic polynomial.} With~\eqref{eq:pd} written as
$C_i(s)=\big[(K_{p,i}+K_{d,i}N)s+K_{p,i}N\big]/(s+N)$ and the plant
numerators $n_\theta(s)=n_2s^2+n_1s+n_0$ and $n_\alpha(s)=m_2s^2$, the
closed-loop characteristic equation of~\eqref{eq:cs1law} is
\begin{equation}\label{eq:char}
D(s)(s+N)+n_\theta(s)\big[(K_{p\theta}+K_{d\theta}N)s+K_{p\theta}N\big]
+n_\alpha(s)\big[(K_{p\alpha}+K_{d\alpha}N)s+K_{p\alpha}N\big]=0.
\end{equation}
This polynomial has degree five, and its $s^5$ and $s^4$ coefficients
($1$ and $N+d_3=71.70$) do not depend on the gains. The sum of the five
closed-loop poles is therefore fixed at $-(N+d_3)$, so the fifth pole is
\emph{not} free:
\begin{equation}
p_5=-(N+d_3)-\big(2(-5.6)-19.2-25.6\big)=-15.7~\mathrm s^{-1}.
\end{equation}
This pole is still faster than the dominant pair.

\textbf{Step 3: solve for the gains.} Matching the $s^3,s^2,s^1,s^0$
coefficients of~\eqref{eq:char} with
$\prod_{k=1}^5(s-p_k)=s^5+71.70s^4+1936.5s^3+24971s^2+162900s+493873$
gives four linear equations in the four gains:
\begin{equation}
\underbrace{\begin{bmatrix}
n_2&n_2N&m_2&m_2N\\
n_1+n_2N&n_1N&m_2N&0\\
n_0+n_1N&n_0N&0&0\\
n_0N&0&0&0\end{bmatrix}}_{M}
\begin{bmatrix}K_{p\theta}\\K_{d\theta}\\K_{p\alpha}\\K_{d\alpha}\end{bmatrix}
=\begin{bmatrix}1936.5-557.5\\24971+9373\\162900+42291\\493873-0\end{bmatrix},
\end{equation}
where the right-hand side is the desired coefficient minus the
corresponding coefficient of $D(s)(s+N)$. $M$ is nonsingular because
$n_0\ne0$ and $m_2\ne0$: the plant has no pole--zero cancellation that
would leave a closed-loop pole uncontrollable by the two loops. Solving with
\texttt{M\textbackslash rhs} in MATLAB (\texttt{params\_CS1A\_candidate.m})
gives
\begin{equation}
K_{p\theta}=\KpTheta,\quad K_{d\theta}=\KdTheta,\quad
K_{p\alpha}=\KpAlpha,\quad K_{d\alpha}=\KdAlpha.
\end{equation}
\begin{itemize}[nosep]
\item \textbf{Arm gains are negative.} The RHP zero makes $n_0<0$, so
  the steady-state arm gain is negative. To move the arm towards $+r$, the
  arm loop must command a negative voltage first, tilting the pendulum
  forward.
\item \textbf{Pendulum loop dominates.} $K_{p\alpha}\gg|K_{p\theta}|$,
  as expected: the pendulum loop provides stability, and the arm loop is
  a slow outer correction.
\end{itemize}

\textbf{Step 4: reference prefilter.} The arm controller adds a closed-loop
zero at
$z_\theta=-K_{p\theta}N/(K_{p\theta}+K_{d\theta}N)=-2.23$~s$^{-1}$. This
is slower than the dominant poles and would cause large overshoot. The
prefilter cancels this zero from the reference path and adds a smooth
3~rad/s lag:
\begin{equation}
F(s)=\frac{2.23}{s+2.23}\cdot\frac{3}{s+3},\qquad F(0)=1.
\end{equation}
$F(s)$ is outside the feedback loop, so it changes the overshoot and the
voltage demand \emph{without moving the closed-loop poles}. This is the
standard way to reduce overshoot while keeping the designed stability and
disturbance rejection~\cite{franklin_feedback}.

\subsection{Frequency-domain verification}
The design is verified on the loop transfer function broken at the motor
input, $L(s)=C_\theta(s)G_\theta(s)+C_\alpha(s)G_\alpha(s)$, so that
$V_m=-L(s)V_m$ around the loop. This is the single point that both loops
share, and where a motor-gain error or an actuator delay enters.
$L(s)$ was built in state space (\texttt{ss([C\_theta\_s, C\_alpha\_s])*G})
because the product of high-order transfer-function polynomials is
numerically ill-conditioned.

\fig{cs1a_bode}{Bode diagram of the CS1 loop $L(j\omega)$ at the motor
input. The gain crosses 0~dB at $\omega_c=\WcpcsA$~rad/s with a phase
margin of $\PMcsA^\circ$.}{fig:cs1bode}

Figure~\ref{fig:cs1bode} shows the loop crossing 0~dB at
$\omega_c=\WcpcsA$~rad/s with a phase margin of $\PMcsA^\circ$. The phase
starts near $-270^\circ$ and rises through $-180^\circ$ at about 7~rad/s,
which is characteristic of a loop that must stabilise an unstable pole:
\begin{itemize}[nosep]
\item By the Nyquist criterion, $L$ has one RHP pole ($P=1$), so stability
  requires one anticlockwise encirclement of $-1$. The phase must cross
  $-180^\circ$ from below while $|L|>1$.
\item The resulting gain margin is a band:
  \begin{itemize}[nosep]
  \item the loop gain may fall to $\GMlowcsA\times$ nominal before the
    pendulum is lost;
  \item it may rise without limit in the continuous model.
  \end{itemize}
\item In the implemented discrete loop (Tustin controller, ZOH plant,
  $T_s=1$~ms), the band is $\GMlowcsAz$--$\GMupcsAz\times$, with a phase
  margin of $\PMcsAz^\circ$ and a delay margin of $\DMcsAz$ samples
  ($\DMcsAz$~ms).
\end{itemize}
The lower limit is the practical constraint. Friction at the small
balancing voltages reduces the effective motor gain, and the design
tolerates a reduction of almost half.

\fig[0.75]{cs1a_rlocus}{Root locus of $kL(s)$ as the loop-gain multiplier
$k$ varies. Stars mark the designed closed-loop poles at $k=1$. The
diamonds mark the poles at $k=\GMlowcsA$, where a complex pair is on the
imaginary axis.}{fig:cs1rl}

Figure~\ref{fig:cs1rl} gives the same information as a root locus:
\begin{itemize}[nosep]
\item As $k$ increases from zero, the branches from the unstable pole
  $s=+9.95$ and the integrator pole $s=0$ meet on the real axis and break
  away as a complex pair. This pair crosses into the left half-plane at
  $s\approx\pm7.3j$ when $k=\GMlowcsA$, which is the lower gain margin and
  the $-180^\circ$ crossing frequency of Figure~\ref{fig:cs1bode}.
\item At $k=1$ the five poles sit exactly at the targets of
  Step~1, confirming the algebra.
\item Near $k=1$, the dominant pair moves only slowly along its branch
  towards the nearby open-loop zeros, so moderate gain errors change the
  damping only slightly.
\end{itemize}

\fig[0.75]{cs1a_sensitivity}{Sensitivity $S(j\omega)$ and complementary
sensitivity $T(j\omega)$ of the CS1 loop.}{fig:cs1sens}

The peak sensitivity is $M_s=\MscsA$ (Figure~\ref{fig:cs1sens}), within
the usual guideline of $M_s<2$~\cite{franklin_feedback}. Disturbances
below about 5~rad/s, such as a bump or friction, are attenuated. The
peak near $\omega_c$ amplifies disturbances by at most $\MscsA\times$.
Above 100~rad/s, $|T|$ rolls off, so encoder quantisation noise is
attenuated before it reaches the motor.

\subsection{Discretisation}
The controllers are mapped to the $z$-domain with the bilinear (Tustin)
transform $s\leftarrow\frac{2}{T_s}\frac{z-1}{z+1}$, using MATLAB
\texttt{c2d(C\_s, ts, 'tustin')}~\cite{mathworks_c2d}:
\begin{align}
C_\alpha(z)&=\frac{181.770-179.271z^{-1}}{1-0.94175z^{-1}},\qquad
C_\theta(z)=\frac{-104.438+104.205z^{-1}}{1-0.94175z^{-1}},\\
F(z)&=\frac{1.668\times10^{-6}(1+2z^{-1}+z^{-2})}{1-1.99478z^{-1}+0.99478z^{-2}}.
\end{align}
The derivative-filter pole maps to $z=0.94175$, close to $e^{-NT_s}=0.94176$.
Tustin was chosen for two reasons:
\begin{itemize}[nosep]
\item it maps the stable left half-plane onto the unit disc, so a stable
  $C(s)$ always gives a stable $C(z)$;
\item it preserves the frequency response up to the frequency-warping
  term, which is negligible here because $\omega_cT_s\approx0.02$.
\end{itemize}
Figure~\ref{fig:cs1c2d} confirms that $C_\alpha(z)$ and $C_\alpha(s)$ are
indistinguishable up to 1000~rad/s. The sampling rate is 50 times the
crossover frequency, and the extra phase lag from the ZOH, about
$\omega_cT_s/2=0.6^\circ$, is negligible.

\fig[0.7]{cs1a_c2d}{Magnitude of the pendulum controller $C_\alpha(s)$
and of its Tustin discretisation $C_\alpha(z)$ at $T_s=1$~ms.}{fig:cs1c2d}

\subsection{Implementation}
The control law~\eqref{eq:cs1law} runs in one MATLAB Function block
(\texttt{ctrl\_CS1A.slx}). That block is referenced by the linear
simulation, the nonlinear simulation and the hardware model, so all three
use the same code. Each stage $b(z^{-1})/a(z^{-1})$ is evaluated as a
difference equation in direct form~II transposed, which needs only $n$
states for an $n$th-order stage. The block also does the following:
\begin{itemize}[nosep]
\item When the controller is disabled (\texttt{reset}=1), it clears every
  stage state and outputs 0~V, so enabling it causes no voltage jump.
\item It saturates $V_m$ to $\pm10$~V.
\end{itemize}
The coefficients come from \texttt{params\_CS1A\_candidate.m}, which
\texttt{Hardware\_settings.m} runs before every build.
\begin{lstlisting}
[r_f, wF]     = lti_step(bF, aF, r, wF);              % F(z)
[u_theta, wT] = lti_step(bT, aT, r_f - theta, wT);    % C_theta(z)
[u_alpha, wA] = lti_step(bA, aA, 0 - alpha, wA);      % C_alpha(z)
u = min(max(u_theta + u_alpha, -u_max), u_max);        % +/-10 V limit

function [y, w] = lti_step(b, a, x, w)   % y = b(z^-1)/a(z^-1) x, DF-II T
n = numel(a) - 1;  y = b(1)*x;
if n > 0
    y = y + w(1);
    for i = 1:n-1, w(i) = b(i+1)*x + w(i+1) - a(i+1)*y; end
    w(n) = b(n+1)*x - a(n+1)*y;
end
end
\end{lstlisting}

\subsection{Simulation results}
\fig{cs1a_sim}{CS1 simulated response to a $45^\circ$ arm step at $t=1$~s
on the linear and nonlinear plant models: arm angle $\theta$, pendulum
angle $\alpha$ and motor voltage $V_m$.}{fig:cs1sim}
\fig[0.8]{cs1a_sim_zoom}{Zoom on the CS1 arm transient of
Figure~\ref{fig:cs1sim}, with the $\pm2\%$ settling band.}{fig:cs1zoom}

Figures~\ref{fig:cs1sim} and~\ref{fig:cs1zoom} show the response:
\begin{itemize}[nosep]
\item \textbf{Initial undershoot.} The arm first moves slightly negative,
  the signature of the RHP zero. The voltage goes negative first to tilt
  the pendulum forward.
\item \textbf{Specifications met.} The arm then rises without overshoot
  ($M_p=\OScsAN\%$) and settles within 2\% in $t_s=\TscsAN$~s on the
  nonlinear plant, well inside both requirements.
\item \textbf{Prefilter effect.} The settling time is longer than the
  0.71~s dominant-pole estimate because the 3~rad/s prefilter lag deliberately
  slows the \emph{reference}. This keeps the peak voltage at only
  $\UcsAN$~V and the pendulum tilt at $\AlcsAN^\circ$.
\item \textbf{Linear and nonlinear agree.} The two curves coincide,
  because $\alpha$ stays small.
\end{itemize}

\subsection{Experimental results}
\fig{hw_cs1a_run04}{CS1 running on the QUBE-Servo~2 (run~04). The
pendulum is caught at the start of the trace and then balanced for
\hwT~s, during which it was bumped by hand several times.}{fig:cs1hw}
\fig[0.8]{hw_cs1a_bump}{Zoom on Figure~\ref{fig:cs1hw} around the
largest bump: pendulum angle (left axis) and motor voltage (right
axis).}{fig:cs1bump}

Figure~\ref{fig:cs1hw} shows CS1 balancing on the hardware for \hwT~s:
\begin{itemize}[nosep]
\item \textbf{Balancing.} Excluding the catch transient, the pendulum
  angle has a standard deviation of $\hwAlphaStd^\circ$, about five
  encoder counts, and an RMS voltage of \hwVrms~V.
\item \textbf{Bump rejection.} The pendulum was bumped by hand several
  times. The largest bump tilted it by $\hwBumpPeak^\circ$, and
  Figure~\ref{fig:cs1bump} shows it back within $\pm2^\circ$ in about
  \hwRecover~s. The voltage briefly reaches about 5~V, and the response is
  well damped, as the $M_s=\MscsA$ of
  Figure~\ref{fig:cs1sens} predicts.
\item \textbf{Voltage chatter.} The persistent $\pm1$~V ripple in $V_m$
  comes from the derivative terms acting on the quantised encoder
  signals.
\item \textbf{Arm drift.} The arm drifts slowly through a range of
  about $\pm110^\circ$. Without integral action, Coulomb friction lets the
  arm rest anywhere inside a dead band, and each bump moves that resting
  point. \todo{Confirm what the operator was doing with the arm and the
  reference during run~04.}
\end{itemize}
\todo{Record a hardware $45^\circ$ arm step with CS1 (with the reference
logged) and add its plot and its measured $M_p$, $t_s$ and steady-state
error here. Runs 01--04 contain balancing and bumps only, with no logged
step.}

%=======================================================================
\section{Control System 2: discrete-time state feedback with an observer}\label{sec:cs2}
\subsection{Control structure and why it was chosen}
State feedback uses all four states, so it can place every closed-loop
pole directly. CS1 could only place five poles through four gains, with
the fifth fixed by the pole-sum constraint. The design is done entirely
in discrete time on the ZOH model $(A,B,C)$ of Section~\ref{sec:model},
so no further approximation is introduced at implementation. Three
additions turn the DP2 SFC into a controller that works on the hardware:
\begin{itemize}[nosep]
\item a \textbf{Luenberger observer}, because only $\theta$ and $\alpha$ are measured;
\item \textbf{integral action} on the arm error, to remove the offset
  caused by friction;
\item a \textbf{first-order reference filter}, to keep a $45^\circ$ step
  inside the voltage limit.
\end{itemize}
The control law is
\begin{align}
r_f[k+1]&=r_f[k]+\tfrac{T_s}{\tau}\big(r-r_f[k]\big),\qquad\tau=0.2~\mathrm s,\\
x_i[k+1]&=x_i[k]+T_s\big(\theta[k]-r\big),\\
u[k]&=\operatorname{sat}_{\pm10}\Big[-F\big(\hat x[k]-x^*\big)-F_ix_i[k]\Big],\qquad x^*=[r_f\ 0\ 0\ 0]^T,\\
\hat x[k+1]&=A\hat x[k]+Bu[k]+L_o\big(y[k]-C\hat x[k]\big).
\end{align}

\subsection{Controllability and observability}
The controllability matrix $\mathcal C=[B\ AB\ A^2B\ A^3B]$ and the
observability matrix $\mathcal O=[C;\,CA;\,CA^2;\,CA^3]$ both have rank
4 (MATLAB \texttt{rank(ctrb(A,B))} and \texttt{rank(obsv(A,C))}). All
closed-loop poles can therefore be placed, and the full state can be
reconstructed from the two encoders. The open-loop discrete poles include
$z=e^{9.95T_s}=1.00999$, outside the unit circle, which confirms the
unstable plant.

\subsection{Pole selection}
The poles are chosen in the $s$-plane from the specifications and mapped
to the $z$-plane with $z=e^{sT_s}$. This mapping is exact for the ZOH
model, since $A=e^{A_cT_s}$:
\begin{itemize}[nosep]
\item \textbf{SFC.} A dominant pair with $\zeta=0.7$ and $\omega_n=8$~rad/s
  ($-5.6\pm5.71j$), plus $-19.2$ and $-25.6$. This is the same
  bandwidth as CS1. The DP2 poles ($\omega_n=2.5$) met the specification
  in simulation, but they left the loop stable only while the motor
  delivered at least $0.87\times$ its modelled gain. The CS1 hardware
  runs showed the real effective gain is lower, so the faster set, which
  tolerates $\GMlowcsB\times$, was adopted.
\item \textbf{Integrator.} The pole is at $-1$~s$^{-1}$. It was first
  placed at $-8$, but that roughly tripled the arm gain, so a small bump
  saturated the voltage and the pendulum fell. At $-1$, the gains stay
  close to plain SFC and friction offsets are removed in about one second.
\item \textbf{Observer.} The poles are at $-90,-95,-100,-105$~s$^{-1}$,
  about $4\times$ faster than the fastest SFC pole, so the estimate
  converges well before the controller acts on it. They are not made
  faster, because a faster observer passes more quantisation noise into
  the velocity estimates.
\end{itemize}

\subsection{Gain calculation}
\begin{itemize}[nosep]
\item \textbf{SFC and integral gains.} $F$ and $F_i$ are computed
  together on the augmented system
  $A_a=\begin{bmatrix}A&0\\T_s[1\ 0\ 0\ 0]&1\end{bmatrix}$,
  $B_a=\begin{bmatrix}B\\0\end{bmatrix}$ with
  \texttt{place(Aa, Ba, exp([ps p\_int]*ts))}, giving
  \[F=\FcsB,\qquad F_i=\FicsB.\]
\item \textbf{Observer gain.} By duality, it is
  $L_o=$\texttt{place(A', C', exp(po\_s*ts))'}.
\end{itemize}
Broken at the plant input, the state-feedback loop has a gain margin band
of $\GMlowcsB$--$\GMupcsB\times$ and a phase margin of $\PMcsB^\circ$.
These are wider than CS1, because state feedback uses the measured
velocities rather than a filtered derivative.

\fig[0.6]{cs2_zplane}{CS2 poles in the $z$-plane, zoomed near $z=1$:
open-loop poles of $A$ (one outside the unit circle), the SFC poles and
the observer poles.}{fig:cs2z}
Figure~\ref{fig:cs2z} shows the mapping $z=e^{sT_s}$ at work:
\begin{itemize}[nosep]
\item the unstable open-loop pole sits just outside the unit circle;
\item the SFC poles move it inside;
\item the faster observer poles lie further from $z=1$, because a larger
  $|s|$ maps to a smaller $|z|$.
\end{itemize}

\subsection{Implementation}
Like CS1, the whole controller is one MATLAB Function block
(\texttt{ctrl\_CS2.slx}) shared by both simulations and the hardware. It
includes three practical details:
\begin{itemize}[nosep]
\item \textbf{Bumpless start.} On enable, the observer starts from the
  measured angles with zero velocities, $r_f$ starts from the measured
  $\theta$, and $x_i=0$, so the controller starts without a voltage jump.
\item \textbf{Conditional integration with anti-windup.} The integrator
  updates only while $|\theta-r|\le2^\circ$, and not while the voltage is
  saturated in the direction that would increase it. Integrating during
  the whole $45^\circ$ move produced 44\% overshoot in simulation.
\item \textbf{Observer driven by the saturated voltage.} The observer uses
  the limited $u$, so it models the input the motor actually receives.
\end{itemize}
\begin{lstlisting}
u_unlimited = -F*(xh - x_ref) - Fi*xi;
u = min(max(u_unlimited, -u_max), u_max);
e = theta - r;  dxi = ts*e;
winding_up = (u_unlimited > u_max && -Fi*dxi > 0) || (u_unlimited < -u_max && -Fi*dxi < 0);
if abs(e) <= int_band && ~winding_up && ~reset, xi = xi + dxi; end
xh = A*xh + B*u + Lo*(y - C*xh);                     % Luenberger observer
\end{lstlisting}

\subsection{Simulation results}
\fig{cs2_sim}{CS2 simulated response to a $45^\circ$ arm step at $t=1$~s
on the linear and nonlinear plant models.}{fig:cs2sim}
\fig[0.8]{cs2_sim_zoom}{Zoom on the CS2 arm transient of
Figure~\ref{fig:cs2sim}, with the $\pm2\%$ settling band.}{fig:cs2zoom}

CS2 settles in $t_s=\TscsBN$~s with $M_p=\OScsBN\%$ on the nonlinear
plant (Figures~\ref{fig:cs2sim} and~\ref{fig:cs2zoom}). This is faster
than CS1, because the reference filter ($\tau=0.2$~s) is lighter than
CS1's prefilter. The cost is a higher peak voltage ($\UcsBN$~V) and a
larger pendulum excursion ($\AlcsBN^\circ$), both still well inside
their limits.

\subsection{Experimental results}
CS2 balances the pendulum on the QUBE. The first hardware version fell
over when bumped, which led to the integrator-pole change described
above. \todo{Add the CS2 hardware run: balance, bump and $45^\circ$ step
with the reference logged, with its plot and measured $M_p$, $t_s$ and
steady-state error. No CS2 hardware data is saved in the repository yet.}

%=======================================================================
\section{Comparison of the two control systems}\label{sec:compare}
\begin{table}[H]\centering
\caption{Performance for a $45^\circ$ arm step. Simulation values are from
\texttt{make\_figures.m}; hardware values are still to be measured.}\label{tab:compare}
\begin{tabular}{lcccc}
\toprule
&\multicolumn{2}{c}{\textbf{CS1 (frequency domain)}}&\multicolumn{2}{c}{\textbf{CS2 (SFC + observer)}}\\
&Nonlinear sim.&Hardware&Nonlinear sim.&Hardware\\
\midrule
Overshoot $M_p$ [\%]&\OScsAN&\todo{}&\OScsBN&\todo{}\\
Settling time $t_s$ [s]&\TscsAN&\todo{}&\TscsBN&\todo{}\\
Peak $|V_m|$ [V]&\UcsAN&\todo{}&\UcsBN&\todo{}\\
Peak $|\alpha|$ [deg]&\AlcsAN&\todo{}&\AlcsBN&\todo{}\\
Gain margin band&$\GMlowcsAz$--$\GMupcsAz$&--&$\GMlowcsB$--$\GMupcsB$&--\\
Phase margin [deg]&\PMcsAz&--&\PMcsB&--\\
\bottomrule
\end{tabular}
\end{table}
\todo{Discuss Table~\ref{tab:compare} once the hardware columns are
filled: which controller tracks better, which rejects bumps better, voltage
chatter, and the friction offset (CS2 has integral action, CS1 does not).}

%=======================================================================
\section{Response to feedback from the final demonstration}\label{sec:feedback}
\todo{List each comment or question from the teaching team at the final
demonstration, what the team changed or tested in response, and the
result with a figure or number. This section is marked as its own
criterion.}

%=======================================================================
\section{Conclusions and future work}\label{sec:conclusion}
Two control systems were designed from the same physically derived
model of the QUBE-Servo~2 and implemented as single MATLAB Function
blocks at 1~kHz.

\textbf{CS1.} Two filtered PD controllers with gains computed from
specification-driven pole targets were verified by Bode, root-locus and
sensitivity analysis. They have a $\PMcsAz^\circ$ phase margin and
tolerate a motor gain down to $\GMlowcsAz\times$ nominal. In simulation,
CS1 meets the requirements with no overshoot and $t_s=\TscsAN$~s. On the
hardware, it balanced the pendulum for minutes and recovered from hand
bumps.

\textbf{CS2.} The discrete SFC with integral action and an observer meets
the requirements in simulation, with $M_p=\OScsBN\%$ and
$t_s=\TscsBN$~s. Its integral action removes the friction offset that CS1
leaves.

\textbf{Key lessons.}
\begin{itemize}[nosep]
\item \textbf{Robustness as well as nominal specifications.} Designs that
  met the specifications on the nominal model failed on the hardware or
  under modelled quantisation. Bandwidth had to be chosen for robustness
  to the real motor gain and the encoder resolution.
\item \textbf{One shared code path.} Keeping the simulated and
  implemented code identical made hardware debugging a matter of
  physics, not code.
\end{itemize}

\textbf{Future work.}
\begin{enumerate}[nosep]
\item Identify the motor gain and Coulomb friction of this QUBE
  experimentally and add friction compensation.
\item Add integral action to CS1.
\item Design CS1 directly in the $z$-domain and compare it with the
  $s$-domain design.
\item Add an energy-based swing-up~\cite{astrom2000} so that the system
  starts from hanging down without manual help.
\end{enumerate}
\todo{Revise once hardware step results and demonstration feedback are
added.}

%=======================================================================
\begin{thebibliography}{99}
\bibitem{aguilera2026}
R. P. Aguilera, \emph{41277 Control Design Lecture Notes}, University of
Technology Sydney, Australia, 2026.

\bibitem{furuta1992}
K. Furuta, M. Yamakita and S. Kobayashi, ``Swing-up control of inverted
pendulum using pseudo-state feedback,'' \emph{Proc. Inst. Mech. Eng. Part
I: J. Syst. Control Eng.}, vol.~206, no.~4, pp.~263--269, 1992.

\bibitem{astrom2000}
K. J. Åström and K. Furuta, ``Swinging up a pendulum by energy control,''
\emph{Automatica}, vol.~36, no.~2, pp.~287--295, 2000.

\bibitem{franklin_feedback}
G. F. Franklin, J. D. Powell and A. Emami-Naeini, \emph{Feedback Control of
Dynamic Systems}, 8th ed., Pearson, 2019.

\bibitem{franklin_digital}
G. F. Franklin, J. D. Powell and M. L. Workman, \emph{Digital Control of
Dynamic Systems}, 3rd ed., Addison-Wesley, 1998.

\bibitem{ogata}
K. Ogata, \emph{Modern Control Engineering}, 5th ed., Prentice Hall, 2010.

\bibitem{ctms}
B. Messner and D. Tilbury, \emph{Control Tutorials for MATLAB and
Simulink}, University of Michigan. [Online]. Available:
\url{https://ctms.engin.umich.edu/CTMS}

\bibitem{mathworks_c2d}
The MathWorks, Inc., ``c2d: Convert model from continuous to discrete
time,'' MATLAB documentation. [Online]. Available:
\url{https://www.mathworks.com/help/control/ref/dynamicsystem.c2d.html}

\bibitem{qube_manual}
Quanser Inc., \emph{QUBE-Servo 2 User Manual}, Markham, ON, Canada.

\bibitem{qube_courseware}
Quanser Inc., \emph{QUBE-Servo Workbook -- Student}, v2.0,
\emph{Pendulum State-Space Model}, in-lab exercise 6.

\bibitem{hardware_template}
R. P. Aguilera, \emph{ControLab Rotary Inverted Pendulum Template},
\texttt{CS01\_RIP\_QS2\_T41\_r2025b}, University of Technology Sydney, 2026.
\end{thebibliography}

\end{document}
